\begin{longtable}{TD}
\caption{Notation used throughout the paper.}
\label{tab:notation}\\

\toprule
Symbol & Meaning \\
\midrule
\endfirsthead

\toprule
Symbol & Meaning \\
\midrule
\endhead

\bottomrule
\endfoot

\termrow{$\mathcal{K}$}{
The set of task kinds.
}

\termrow{$\mathcal{I}$}{
The set of task instances.
}

\termrow{$\mathcal{A}$}{
The set of task attempts.
}

\termrow{$i \in \mathcal{I}$}{
A task instance.
}

\termrow{$a \in \mathcal{A}$}{
A task attempt.
}

\termrow{$k(i) \in \mathcal{K}$}{
The task kind of task instance $i$.
}

\termrow{$i(a) \in \mathcal{I}$}{
The task instance executed by attempt $a$.
}

\termrow{$k_a$}{
A shorthand for $k(i(a))$, the task kind of attempt $a$.
}

\termrow{$\mathcal{D}$}{
The set of influence domains.
}

\termrow{$\mathcal{F}$}{
The set of canonical features.
}

\termrow{$z_a$}{
The canonical feature vector available for attempt $a$ after feature mapping
and normalization.
}

\termrow{$\mathbf{x}_a$}{
The intrinsic-feature projection of $z_a$, containing task-instance features
such as payload size, record count, batch size, operation mode, or other
domain-specific work descriptors.
}

\termrow{$\mathbf{c}_a$}{
The execution-context projection of $z_a$, containing context features such as
worker state, node pressure, dependency health, network condition, disk
pressure, cache state, retry state, tenant pressure, or domain-specific state.
}

\termrow{$\mathcal{R}$}{
The set of active resource-demand dimensions.
}

\termrow{$\mathcal{O}$}{
The set of worker-occupancy dimensions.
}

\termrow{$\mathcal{Q}$}{
The set of cost dimensions, defined as $\mathcal{Q} = \mathcal{R} \cup
\mathcal{O}$.
}

\termrow{$q \in \mathcal{Q}$}{
A cost dimension, either a resource-demand dimension or a worker-occupancy
dimension.
}

\termrow{$\widehat{D}_r(a)$}{
Estimated resource demand for resource dimension $r \in \mathcal{R}$ during
attempt $a$.
}

\termrow{$\widehat{U}_o(a)$}{
Estimated worker occupancy for occupancy dimension $o \in \mathcal{O}$ during
attempt $a$.
}

\termrow{$\widehat{\mathbf{D}}(a)$}{
The resource-demand estimate vector for attempt $a$.
}

\termrow{$\widehat{\mathbf{U}}(a)$}{
The worker-occupancy estimate vector for attempt $a$.
}

\termrow{$\widehat{\mathbf{C}}(a)$}{
The structured attempt-level cost estimate, represented as
$(\widehat{\mathbf{D}}(a), \widehat{\mathbf{U}}(a))$.
}

\termrow{$\widehat{Y}_q(a)$}{
The estimated value for cost dimension $q \in \mathcal{Q}$ during attempt $a$.
This is a dimension-generic notation used when the same expression applies to
both resource-demand and worker-occupancy dimensions.
}

\termrow{$Y_q(a)$}{
The observed attempt-attributed value for cost dimension $q$, when such an
observation is available. For $q \in \mathcal{R}$ this is observed resource
usage; for $q \in \mathcal{O}$ this is observed worker occupancy.
}

\termrow{$\widehat{B}_q$}{
The estimated intrinsic baseline function for cost dimension $q$.
}

\termrow{$\widehat{A}_q$}{
The estimated contextual amplification factor for cost dimension $q$.
}

\termrow{$\widehat{O}_q$}{
The estimated additive overhead function for cost dimension $q$.
}

\termrow{$\epsilon_q(a)$}{
Residual error for cost dimension $q$, defined when $Y_q(a)$ is available as
$Y_q(a) - \widehat{Y}_q(a)$.
}

\termrow{$\widehat{Y}_{q,p}(a)$}{
The estimated $p$-quantile or percentile-level value for cost dimension $q$
during attempt $a$.
}

\termrow{$\mathrm{Risk}(a)$}{
The operational-risk estimate for attempt $a$.
}

\termrow{$\mathcal{G}_k$}{
The Task Behavior Graph associated with task kind $k$.
}

\termrow{$V_k$}{
The set of nodes in Task Behavior Graph $\mathcal{G}_k$.
}

\termrow{$E_k$}{
The set of edges in Task Behavior Graph $\mathcal{G}_k$.
}

\termrow{$v \in V_k$}{
A node in Task Behavior Graph $\mathcal{G}_k$.
}

\termrow{$e \in E_k$}{
An edge in Task Behavior Graph $\mathcal{G}_k$.
}

\termrow{$\delta(v)$}{
The influence domain assigned to graph node $v$.
}

\termrow{$w_e$}{
The estimated sensitivity or weight associated with graph edge $e$.
}

\termrow{$\gamma_e$}{
The confidence value associated with graph edge $e$.
}

\end{longtable}
